\section{Dynkin's formula, every bidegree, and the polynomial through degree six}\label{sec:dynkin}
\subsection{Dynkin's complete commutator formula}
The structural theorem becomes completely explicit by applying the Dynkin
operator to the associative logarithm. The classical coefficient formula is
due to Dynkin \cite{dynkin}; modern computational treatments include
\cite{casasmurua}. The proof below derives its coefficients directly.

\begin{theorem}[Dynkin's all-terms formula]\label{thm:dynkin}
With $d=\sum_{i=1}^k(r_i+s_i)$,
\begin{equation}\label{eq:dynkin}
 \boxed{\displaystyle
 Z(X,Y)=\sum_{k\geq1}\frac{(-1)^{k-1}}{k}
 \sum_{\substack{r_i,s_i\geq0\\r_i+s_i>0\ (1\leq i\leq k)}}
 \frac{R\bigl(X^{r_1}Y^{s_1}\cdots X^{r_k}Y^{s_k}\bigr)}
      {d\,\prod_{i=1}^k r_i!s_i!}.}
\end{equation}
Equivalently, in terms of the word coefficients \eqref{eq:wordcut},
\begin{equation}\label{eq:dynkinword}
 Z_n=\frac1n\sum_{|w|=n}c(w)\,R(w),
\end{equation}
and, separating the bidegrees,
\begin{equation}\label{eq:bidegreeproj}
 \boxed{\displaystyle
 Z_{p,q}=\frac1{p+q}
 \sum_{\substack{w\in\{X,Y\}^{p+q}\\\#_Xw=p,\ \#_Yw=q}}
 c(w)\,R(w),\qquad p+q\geq1.}
\end{equation}
In each total degree the sums are finite. A summand of \eqref{eq:dynkin}
vanishes if $s_k>1$, or if $s_k=0$ and $r_k>1$.
\end{theorem}
\begin{proof}
By \cref{thm:formalbch}, each $Z_n$ is a homogeneous primitive element of
degree $n$. \Cref{lem:dsw} therefore gives $Z_n=\frac1nD(Z_n)$. Since $D$ is
linear and $D(w)=R(w)$ on words, applying $D/n$ to the word expansion
$Z_n=\sum_{|w|=n}c(w)w$ of \eqref{eq:homblocks} gives \eqref{eq:dynkinword}.
Applying it instead to the block form of \eqref{eq:homblocks} replaces each
block word $X^{r_1}Y^{s_1}\cdots X^{r_k}Y^{s_k}$ by its right-nested bracket
and divides by its length $d=n$; summing over $n$ gives \eqref{eq:dynkin}.
Finally, $Z_n=\sum_{p+q=n}Z_{p,q}$ and the operator $D$ preserves bidegree
(it maps a word to a combination of words with the same letters), so the
bidegree-$(p,q)$ part of \eqref{eq:dynkinword} is \eqref{eq:bidegreeproj}.
If $s_k>1$ the two innermost letters of the bracketed word are $Y,Y$, and
if $s_k=0$, $r_k>1$ they are $X,X$; in either case the innermost bracket
$[a_{d-1},a_d]$ vanishes and so does the whole nested bracket.
\end{proof}
The two vanishing tests are sufficient but not necessary: a run of equal
letters may continue across a block boundary, for instance for the three
blocks $X$, $XY$, $Y$, whose concatenation $XXYY$ has innermost bracket
$[Y,Y]=0$ although $s_k=1$. Notice also that one must project a
\emph{whole homogeneous Lie polynomial}; applying $D/n$ to an arbitrary
associative polynomial generally changes it. The right-nested brackets form a
spanning family of $\cL_n$, not a basis, so different displays can represent
the same Lie element by Jacobi identities. This is redundancy in the
representation, not ambiguity in $Z$. Formula \eqref{eq:bidegreeproj},
combined with the finite cut formula \eqref{eq:wordcut}, is a particularly
compact answer to the request for \emph{all} terms: there are no unspecified
coefficients or auxiliary recurrences.

\subsection{Symmetries, associativity, and functoriality}
\begin{proposition}\label{prop:symmetry}
As universal formal series, with $T$ a third letter,
\begin{align}
 Z(Y,X)&=-Z(-X,-Y),\label{eq:swap}\\
 Z_n(Y,X)&=(-1)^{n+1}Z_n(X,Y),\label{eq:parity}\\
 \rev\bigl(Z_n(X,Y)\bigr)&=(-1)^{n+1}Z_n(X,Y),
 \qquad\text{that is,}\qquad c(\rev w)=(-1)^{|w|+1}c(w),\label{eq:reversal}\\
 Z\bigl(Z(X,Y),T\bigr)&=Z\bigl(X,Z(Y,T)\bigr),\label{eq:assoc}\\
 Z(X,0)&=Z(0,X)=X,\qquad Z(X,-X)=0.\label{eq:unitinverse}
\end{align}
Here $\rev$ reverses every word. Moreover, every Lie-algebra homomorphism
$f:\g\to\h$ satisfies $f(Z_n(x,y))=Z_n(fx,fy)$ for all $n$, and hence
commutes with the BCH series wherever formal substitution or analytic
convergence is valid.
\end{proposition}
\begin{proof}
Taking inverses in $e^{-X}e^{-Y}=e^{Z(-X,-Y)}$ gives
$e^Ye^X=(e^{Z(-X,-Y)})^{-1}=e^{-Z(-X,-Y)}$, using $(e^A)^{-1}=e^{-A}$ from
\cref{lem:explog}. Since also $e^Ye^X=e^{Z(Y,X)}$, uniqueness of the formal
logarithm gives \eqref{eq:swap}. Replacing $(X,Y)$ by $(-X,-Y)$ multiplies
the degree-$n$ part by $(-1)^n$, so comparing homogeneous parts in
\eqref{eq:swap} gives \eqref{eq:parity}. Word reversal $\rev$ is an
antiautomorphism of $\widehat T$: $\rev(uv)=\rev(v)\rev(u)$. It fixes $X$
and $Y$, hence $\rev(e^X)=e^X$, $\rev(e^Y)=e^Y$, and
$\rev(e^Xe^Y)=e^Ye^X$. Because $\rev$ is linear, continuous, and
multiplicative on powers of a single element, it commutes with the
logarithm series: $\rev(\log(1+U))=\log(1+\rev U)$. Therefore
$\rev(Z(X,Y))=\log(e^Ye^X)=Z(Y,X)$, and \eqref{eq:parity} gives
\eqref{eq:reversal}; the coefficient statement is its coordinate form.

For \eqref{eq:assoc}, work in the completion of the free algebra on three
letters $X,Y,T$. Both sides exponentiate to $e^Xe^Ye^T$: the left side
because $e^{Z(Z(X,Y),T)}=e^{Z(X,Y)}e^T=e^Xe^Ye^T$, the right side
similarly; here we use that $Z(A,B)$ for positive-degree $A,B$ is defined
by substitution into the universal series and satisfies $e^{Z(A,B)}=e^Ae^B$
by the substitution homomorphism. Uniqueness of the logarithm proves
\eqref{eq:assoc}. The identities \eqref{eq:unitinverse} follow in the same
way from $e^Xe^0=e^X$ and $e^Xe^{-X}=1$.

A Lie homomorphism $f$ preserves every iterated bracket, hence every
right-nested bracket $R(w)$ evaluated at $(x,y)$, hence every $Z_n$ by
\eqref{eq:dynkinword}. In a convergent analytic setting a continuous linear
map preserves the limit of the series; a linear map between
finite-dimensional spaces is automatically continuous.
\end{proof}
Both the swap symmetry and the reversal symmetry provide useful independent
checks on large coefficient tables, and they are used as such in
\cref{sec:computation}.

\subsection{Every term through degree six}
With the right-nesting convention \eqref{eq:rightbracket}, the terms printed
on the source page are exactly
\begin{align}
 Z_1={}&X+Y,\label{eq:z1}\\
 Z_2={}&\tfrac12\Rb{XY}=\tfrac12[X,Y],\label{eq:z2}\\
 Z_3={}&\tfrac1{12}\bigl(\Rb{XXY}+\Rb{YYX}\bigr)
       =\tfrac1{12}[X,[X,Y]]+\tfrac1{12}[Y,[Y,X]],\label{eq:z3}\\
 Z_4={}&-\tfrac1{24}\Rb{YXXY}=-\tfrac1{24}[Y,[X,[X,Y]]],\label{eq:z4}\\
 Z_5={}&-\tfrac1{720}\bigl(\Rb{YYYYX}+\Rb{XXXXY}\bigr)
       +\tfrac1{360}\bigl(\Rb{XYYYX}+\Rb{YXXXY}\bigr)
       \notag\\
     &+\tfrac1{120}\bigl(\Rb{YXYXY}+\Rb{XYXYX}\bigr),\label{eq:z5}\\
 Z_6={}&\tfrac1{240}\Rb{XYXYXY}
       +\tfrac1{720}\bigl(\Rb{XYXXXY}-\Rb{XXYYXY}\bigr)
       \notag\\
     &+\tfrac1{1440}\bigl(\Rb{XYYYXY}-\Rb{XXYXXY}\bigr).\label{eq:z6}
\end{align}
For example $\Rb{XYXXXY}=[X,[Y,[X,[X,[X,Y]]]]]$. There is no omitted term
of any of these degrees; equations \eqref{eq:wordcut} and
\eqref{eq:dynkin} define every subsequent degree.

For direct comparison with associative matrix multiplication, the first
four degrees are
\begin{align}
 Z_2={}&\tfrac12(XY-YX),\label{eq:assoc2}\\
 Z_3={}&\tfrac1{12}\bigl(X^2Y+XY^2-2XYX+Y^2X+YX^2-2YXY\bigr),\label{eq:assoc3}\\
 Z_4={}&\tfrac1{24}\bigl(X^2Y^2-2XYXY-Y^2X^2+2YXYX\bigr).\label{eq:assoc4}
\end{align}

To turn displayed brackets into word coefficients we use the following
explicit expansion. For $w=a_1\cdots a_n$ and a subset
$S\subseteq\{1,\ldots,n-1\}$, let $w_S$ be the subword of $a_1\cdots a_{n-1}$
at the positions in $S$ and $w_{S^c}$ the subword at the remaining positions
of $\{1,\ldots,n-1\}$; in contrast with \eqref{eq:unshuffle}, the index set
here omits the last position. Then
\begin{equation}\label{eq:bracketexpand}
 R(w)=\sum_{S\subseteq\{1,\ldots,n-1\}}(-1)^{n-1-|S|}\,w_S\,a_n\,\rev(w_{S^c}).
\end{equation}
\begin{proof}[Proof of \eqref{eq:bracketexpand}]
Induct on $n$; the case $n=1$ reads $R(a_1)=a_1$. For $n\geq2$,
$R(w)=a_1R(a_2\cdots a_n)-R(a_2\cdots a_n)a_1$. By the induction hypothesis
each term of $R(a_2\cdots a_n)$ has the form
$\pm(\text{subword of }a_2\cdots a_{n-1})\,a_n\,(\text{reversed complementary
subword})$. Placing $a_1$ on the left adds position $1$ to $S$ and keeps the
sign; placing $a_1$ on the right, with a minus sign, appends $a_1$ at the end
of the reversed complement, which is exactly where position $1$ belongs in
$\rev(w_{S^c})$ when $1\notin S$, and changes the sign, matching the factor
$(-1)^{n-1-|S|}$ because $|S^c|$ has increased by one.
\end{proof}

\begin{proposition}\label{prop:lowdegree}
Equations \eqref{eq:z1}--\eqref{eq:z6} are the homogeneous components of
degrees one through six of $Z(X,Y)$, and \eqref{eq:assoc2}--\eqref{eq:assoc4}
are their associative expansions.
\end{proposition}
\begin{proof}
Degree one is \cref{thm:formalbch}. For degree two, \eqref{eq:wordcut}
gives $c(XY)=\tfrac12$ and $c(YX)=-\tfrac12$, while
$c(XX)=a(XX)-\tfrac12a(X)a(X)=\tfrac12-\tfrac12=0$ and likewise
$c(YY)=0$. Hence $Z_2=\tfrac12(XY-YX)=\tfrac12[X,Y]$, which is
\eqref{eq:z2} and \eqref{eq:assoc2}.

For degree three the cut formula gives the following values. For $XXY$ we
computed $\tfrac1{12}$; by reversal symmetry \eqref{eq:reversal} with the
sign $(-1)^{4}=+1$, $c(YXX)=\tfrac1{12}$, and by swap symmetry
\eqref{eq:parity}, $c(YYX)=c(XYY)=\tfrac1{12}$. For $XYX$ we computed
$-\tfrac16$, and by swap, $c(YXY)=-\tfrac16$. For $XXX$: the one-piece
contribution is $\tfrac16$, the two cuts contribute
$-\tfrac12(\tfrac12+\tfrac12)=-\tfrac12$, and the three-piece cut
contributes $\tfrac13$; the sum is zero, and likewise $c(YYY)=0$. Thus
\[
 Z_3=\tfrac1{12}\bigl(XXY+YXX+YYX+XYY-2XYX-2YXY\bigr),
\]
which is \eqref{eq:assoc3}. On the other hand, expanding the brackets,
\[
 [X,[X,Y]]=X^2Y-2XYX+YX^2,\qquad
 [Y,[Y,X]]=Y^2X-2YXY+XY^2,
\]
and $\tfrac1{12}$ times their sum is the same polynomial. This proves
\eqref{eq:z3}.

For degree four, direct evaluation of \eqref{eq:wordcut} on the sixteen
words gives nonzero values only for
\[
 c(XXYY)=\tfrac1{24},\quad c(XYXY)=-\tfrac1{12},\quad
 c(YXYX)=\tfrac1{12},\quad c(YYXX)=-\tfrac1{24},
\]
and zero on the remaining twelve words. We verify one representative of
each type. For $XXYY$: one piece gives $a(XXYY)=\tfrac14$; the three
two-piece cuts $X\mid XYY,\ XX\mid YY,\ XXY\mid Y$ give
$-\tfrac12(\tfrac12+\tfrac14+\tfrac12)=-\tfrac58$; the three cuts into three pieces,
$X\mid X\mid YY,\ X\mid XY\mid Y,\ XX\mid Y\mid Y$, give
$\tfrac13(\tfrac12+1+\tfrac12)=\tfrac23$; the four-piece cut gives
$-\tfrac14$; the sum is $\tfrac14-\tfrac58+\tfrac23-\tfrac14=\tfrac1{24}$.
For $XYXY$: one piece gives $0$; two-piece cuts $X\mid YXY$ ($0$),
$XY\mid XY$ ($1$), $XYX\mid Y$ ($0$) give $-\tfrac12$; three-piece cuts
$X\mid Y\mid XY$, $X\mid YX\mid Y$, $XY\mid X\mid Y$ give
$\tfrac13(1+0+1)=\tfrac23$; four pieces give $-\tfrac14$; the sum is
$-\tfrac12+\tfrac23-\tfrac14=-\tfrac1{12}$. The values at $YYXX$ and $YXYX$
follow by reversal symmetry with sign $(-1)^5=-1$. For a word with zero
coefficient, take $XXXY$: one piece gives $a(XXXY)=\tfrac16$; the two-piece
cuts $X\mid XXY$, $XX\mid XY$, $XXX\mid Y$ give
$-\tfrac12(\tfrac12+\tfrac12+\tfrac16)=-\tfrac7{12}$; the three-piece cuts
$X\mid X\mid XY$ ($1$), $X\mid XX\mid Y$ ($\tfrac12$), $XX\mid X\mid Y$
($\tfrac12$) give $\tfrac13\cdot2=\tfrac23$; four pieces give $-\tfrac14$;
the sum is $\tfrac16-\tfrac7{12}+\tfrac23-\tfrac14=0$. The
remaining zero coefficients are checked in the same way; the eleven other
words are $XXXX$, $YYYY$, $XXYX$, $XYXX$, $YXXX$, $XYYY$, $YXYY$, $YYXY$,
$YYYX$, $XYYX$, $YXXY$, and each cut sum vanishes (for the last two, for
example, the contributions are $0-\tfrac12(0+0+\tfrac12)+\tfrac13(0+\tfrac12+1)
-\tfrac14=0$ and $0-\tfrac12(\tfrac12+0+0)+\tfrac13(1+\tfrac12+0)-\tfrac14=0$,
the cuts being listed from left to right by the position of the first break).
Thus \eqref{eq:assoc4} holds. Expanding the bracket,
\[
 [Y,[X,[X,Y]]]=[Y,\,X^2Y-2XYX+YX^2]
 =2XYXY-2YXYX+Y^2X^2-X^2Y^2,
\]
and $-\tfrac1{24}$ times this is \eqref{eq:assoc4}. This proves
\eqref{eq:z4}.

For degrees five and six we give a finite coefficientwise certificate.
Expand each right-nested bracket in \eqref{eq:z5} and \eqref{eq:z6} by
\eqref{eq:bracketexpand} and collect the coefficient of every word $w$ of
length five or six. Independently, evaluate the finite cut sum
\eqref{eq:wordcut} for every such word. \Cref{app:tables} prints all nonzero
values of $c(w)$ through degree six (there are $30$ nonzero coefficients in
degree five and $28$ in degree six), and both computations give zero on
every word absent from the table. Thus the two polynomials agree in every
one of the $32$ and $64$ coordinates respectively. An equality of
noncommutative polynomials is fully determined by these word coefficients,
so the displayed formulas are identities in the free associative algebra,
hence in every associative specialization. Only the cut-formula side is
printed in \cref{app:tables}; the bracket-expansion side is reproduced,
coefficient by coefficient including the zeros, by the table generator and
by the verifiers described in \cref{sec:computation}, which additionally
recompute both sides from the differential recursions of
\cref{sec:integral}.
\end{proof}
These displays do not select a linearly independent basis of the free Lie
algebra; for instance $[X,[Y,[X,Y]]]=[Y,[X,[X,Y]]]$, since the difference is
$[[X,Y],[X,Y]]=0$ by \eqref{eq:derivation}. The coefficient of $X^4Y$ in
$Z_5$ is $-\tfrac1{720}=B_4^+/4!$, and that of $X^5Y$ in $Z_6$ is
$0=B_5^+/5!$, as the single-$Y$ Bernoulli coefficients of
\cref{cor:linearY} predict.
